\makeabstract
    {
    Is the Universe Symmetric? Reinterpreting Data to Probe Local Lorentz Invariance
    }
    {
    August Castelli\textsuperscript{1}, Daniel Murphy\textsuperscript{1}, Nathan Clayburn\textsuperscript{2}, Andrew Glassford\textsuperscript{1}, Ashley Kim\textsuperscript{1}, Larry Hunter\textsuperscript{1}
    }
    {
    \textsuperscript{1} Department of Physics \& Astronomy, Amherst College, Amherst, MA\\
\textsuperscript{2} Department of Physics, Keene State College, Keene, NH
    }
    {
    We reinterpreted data from an experiment optimized to search for long range spin-dependent interactions to probe Local Lorentz Invariance (LLI). LLI is a principle in physics which states that the laws of nature are the same for all inertial reference frames. LLI is foundational to most modern physics and is violated in many candidates for theories of quantum gravity. A Lorentz symmetry violation could also help explain currently unexplained phenomena, such as matter{\textendash}antimatter asymmetry and dark energy.
    }
    {
    In the original search for long-range spin-dependent interactions, we compared comagnetometer data across different orientations of the internal magnetic field relative to Earth. To test LLI, we instead compared the data across field orientations relative to the fixed stars, allowing us to probe the spatial isotropy, or directional uniformity, of spin in the universe. Each data point carries a timestamp and a magnetic field direction, so we first converted each timestamp to local sidereal time, which is equivalent to right ascension, and then offset the resulting values to account for the direction the field was facing. Our experiment is prone to drifts from systematic effects, which we corrected with a linear regression. Most of our data runs were shorter than 24 hours, meaning any potential Lorentz-violation effect on those runs could also have been erroneously eliminated by the linear regression. Out of our 29 runs, run 26 retained significant drift even after correction, while runs 10, 22, and 28 showed lesser drift that could be mostly corrected by the linear regression. A violation of spatial isotropy should show in our data as a sinusoidal variation with a period of 24 hours. We fit the data to a sine curve to get an amplitude and accompanying 2\ensuremath{\sigma} uncertainty. We also converted the amplitude to equatorial neutron and proton couplings.

    }
    {
    Our final amplitude, with only run 26 excluded, was 20 \ensuremath{\pm} 35 nHz, a null result. Our neutron energy bound was 7.6 \ensuremath{\times} 10{\textasciicircum}-31 eV and our proton energy bound was 7.4 \ensuremath{\times} 10{\textasciicircum}-30 eV.
    }
    {
    Our final results are, at best, roughly a factor of 2 worse than the Hunter lab's prior dedicated LLI experiment. These results can be explained because, unlike most LLI probing-experiments, our runs are relatively short (at most about a day or two), so we are unable to correct for any quadratic drifts. Also, unlike the previous experiment, we are not constantly taking data throughout each run. Taking longer runs would improve our bounds, but a redesign to a dedicated LLI-probing experiment would likely be required for a more significant improvement.

    }

    \newpage
    